\documentclass[11pt]{article}
\usepackage[a4paper,margin=1in]{geometry}
\usepackage{amsmath,amssymb,amsthm,booktabs}
\newtheorem{theorem}{Theorem}
\newtheorem{proposition}[theorem]{Proposition}
\newtheorem{lemma}[theorem]{Lemma}
\newcommand{\Q}{\mathbf Q}
\title{A streamlined computer-assisted proof of\\
Sharipov's second cuboid conjecture}
\author{Computational research draft}
\date{August 11, 2026}

\begin{document}
\maketitle

\begin{abstract}
We prove that Sharipov's second cuboid polynomial is irreducible for every
pair of coprime positive integers $p\ne q$.  The proof reduces factorization
to the $1+4$, $2+3$, and symmetric $5+5$ cases.  We compress the completed
$2+3$ argument into one lossless reduction to
\[
 r^2=4u^8-8u^2+5
\]
and a square-abscissa theorem on
$R^2=4a^5-8a^2+5a$.  The latter uses an exact genus-two 2-descent, an
index-independent Mordell--Weil sieve at $7$ and $23$, and a $7$-adic
Coleman bound.  All denominator and exceptional fibres are retained in
exact certificates.  The proof uses neither a height-bounded search, a
finite congruence sieve as a completeness premise, nor the open direct
normalization claim for the older plane model.
\end{abstract}

\section{Factorization reduction}

For $s=(p/q)^2$, put
\[
\begin{aligned}
P_s(x)={}&x^5+(2+s)(3-2s)x^4
 +(1+10s+4s^2-14s^3+s^4)x^3\\
&-s(1-14s+4s^2+10s^3+s^4)x^2\\
&-s^3(1+2s)(-2+3s)x-s^5
\end{aligned}
\]
and
\[
 Q_{p,q}(t)=q^{20}P_s(t^2/q^4).
\]
The admissible domain is $p,q>0$, $\gcd(p,q)=1$, and $p\ne q$; hence
$s>0$ and $s\ne1$.

\begin{proposition}[Complete factorization reduction]\label{prop:factor}
It is enough to exclude a rational linear factor of $P_s$, a rational
quadratic factor of $P_s$, and a symmetric factorization
\[
 P_s(u^2)=R(u)(-R(-u))
\]
with $R$ irreducible of degree five.
\end{proposition}

\begin{proof}
The scaling $u=t/q^2$ and the nonzero factor $q^{20}$ preserve
reducibility.  If the monic quintic $P_s$ is reducible, one of its
irreducible factors has degree at most two, giving the first two cases.
If $P_s$ is irreducible and $\alpha$ is a root, Capell's composition
criterion says that $P_s(u^2)$ is reducible precisely when $\alpha$ is a
square in $\Q(\alpha)$.  Writing $\alpha=\beta^2$, one has
$\Q(\beta)=\Q(\alpha)$, so the minimal polynomial $R$ of $\beta$ has degree
five and conjugating $\beta$ to $-\beta$ gives the displayed factorization.
Conversely that factorization is reducible.  Finally, monicity and Gauss's
lemma identify rational and integral factorizations of $Q_{p,q}$.
\end{proof}

\section{The linear--quartic case}

\begin{theorem}[Linear-factor exclusion]\label{thm:linear}
For positive rational $s\ne1$, the quintic $P_s$ has no rational root.
\end{theorem}

\begin{proof}
The exact reciprocal quotient sends a rational root to a rational point on
\[
 C_1:\quad w^2=t^5+21t^4+26t^3+10t^2+5t+1.
\]
A Schaefer--Stoll descent in
$\Q\times\Q[z]/(z^4-2z-1)$, with certified class and unit groups, proves
\[
 \dim_{\mathbf F_2}\operatorname{Sel}_2(J_1/\Q)=2,\qquad
 \operatorname{rank}J_1(\Q)=1,\qquad
 \operatorname{Sha}(J_1/\Q)[2]=0.
\]
At $5$, the six known rational points reduce bijectively to
$C_1(\mathbf F_5)$.  An annihilating differential reduces to
\[
 (1+t)\frac{dt}{2w}.
\]
It is nonzero in five disks.  In the remaining disk at $(-1,0)$, the
parameter $w=5z$ gives the exact Strassmann form
\[
 \frac{G(5z)}{25}
 =z\bigl(\text{\rm unit}+5z^2\mathbf Z_5[[z^2]]\bigr).
\]
Thus
\[
 C_1(\Q)=\{\infty,(-1,0),(0,\pm1),(1,\pm8)\}.
\]
The inverse quotient sends these points only to $s=1$, $s=-1$, or poles.
The remaining denominator $t^2+6t+1$ has discriminant $32$ and no rational
zero.  Hence no admissible $s$ occurs.  The descent, local series, quotient
identities, and lift table are replayed by the certificates in
\texttt{research/linear\_factor\_completion/}.
\end{proof}

\section{A lossless reduction of the quadratic--cubic case}

\begin{lemma}[Signed-square constant]\label{lem:signed}
If $X^2+AX+B$ is an integral monic factor of the quintic in $X=t^2$, then
\[
 B=\sigma d^2,\qquad \sigma\in\{+1,-1\},\quad d\in\mathbf Z.
\]
\end{lemma}

\begin{proof}
For a prime $\ell\mid pq$, put $k=v_\ell(pq)$.  Direct coefficient
valuations give the lower Newton polygon with vertices
\[
 (0,10k),(2,2k),(3,0),(5,0),
\]
including $\ell=2,3$.  The root valuations are
$4k,4k,2k,0,0$, so the product of any two roots has even valuation.
Since $B$ divides $-p^{10}q^{10}$, all valuations of $B$ are even.  The
single sign $\sigma$ is physical and global.
\end{proof}

\begin{theorem}[Consolidated $2+3$ reduction]\label{thm:reduction}
Every admissible rational quadratic factor of $P_s$ determines a rational
solution of
\[
 C_3:\quad r^2=4u^8-8u^2+5,                         \tag{1}
\]
and hence a point
\[
 (a,R)=(u^2,ur)
\quad\hbox{on}\quad
 D_+:\ R^2=4a^5-8a^2+5a,                            \tag{2}
\]
with square abscissa.  Every source, base-locus, denominator-zero, and
exceptional fibre omitted by the generic formulas is explicitly finite and
inadmissible.
\end{theorem}

\begin{proof}
We give the exact chain while suppressing large coordinates from the
statement.  Write $s=z^2$, normalize the constant of the quadratic factor as
$b=\sigma z^4k^2$, and write its linear coefficient as $z^2ky$.
Substitution in the two exact remainder equations and elimination of $y$
gives an irreducible $71$-term curve
\[
 Q_\sigma(z,k)=0.
\]
The penultimate subresultant is $L_1y+L_0$ and recovers
$y=-L_0/L_1$.  The resultant of $Q_\sigma$ with the non-boundary factor of
$L_1$ is
\[
\begin{aligned}
2^{20}z^{24}(z-1)^{12}(z+1)^{12}
 &(z^2-z+1)^4(z^2+z+1)^4(z^2+1)^{12}\\
 &\cdot\Phi_{32}(z)^4,
\end{aligned}
\]
where $\Phi_{32}$ is monic with constant $1$ and
$\Phi_{32}(\pm1)=512$.  The quadratic factors have negative discriminant.
Thus the only rational exceptional values are $z=0,\pm1$, namely
$s=0,1$.  This proves that the compact model and inverse are lossless on
the admissible locus.

After putting $S=z^2$, $K=k^2$ and applying two Cayley transforms, the
simultaneous reciprocal involution has fixed generators
$\mathcal U,\mathcal V$.  Both physical signs give the same irreducible
$29$-term curve $C_8(\mathcal U,\mathcal V)=0$ and require
\[
 \mathcal U\mathcal V,\qquad 1-\mathcal U\mathcal V       \tag{3}
\]
to be nonzero rational squares.  Exact quartic functions give maps in both
directions between the admissible normalization of $C_8$ and
\[
 E_1:\quad y^2=x^3-5x+4.                                \tag{4}
\]
Substitution, linear subresultants, and multiplication-back verify both
round trips.  The factored source and inverse denominator resultants leave
only the fibres in Table~\ref{tab:exceptions}; each is a boundary,
nonrational, sign-incompatible, or already fails (3).  In particular the
negative source chart's potentially lost radial fibre has discriminant
$2^8\cdot5$, and the removable elliptic base point has odd $2$-adic
valuation of $\mathcal U\mathcal V$.  No generic open-set assertion replaces
this audit.

Translate (4) by its rational two-torsion point, obtaining coordinates
$(X,Y)$, and put
\[
L=2X+Y-2,\quad M=X^2+2Y-1,\quad
\kappa=\frac{Y+2}{X},\quad h=X+1.
\]
Exact reduction in the function field gives
\[
 XM=(Y+2)L,\qquad
 1-\mathcal U\mathcal V
 =2(X+1)LMg^2                                      \tag{5}
\]
for a recorded rational function $g$.  The factor $2$ in (5) is essential.
Away from the separately audited zeros and poles, (3) implies
\[
 \xi^2=2\kappa h,\qquad
 h^2-(\kappa^2+2)h+\kappa^2+4\kappa-4=0.           \tag{6}
\]

For every odd prime, the Newton polygon of the second equation and the first
equation in (6) show that $v(\kappa)$ is even.  At $2$, the discriminant
$\kappa^4-16\kappa+20$ excludes $v_2(\kappa)=0$ and
$v_2(\kappa)\ge2$; negative valuation has root valuations $0$ and
$2v_2(\kappa)$.  Hence
\[
 \kappa=2\epsilon u^2,\qquad h=\epsilon v^2,\qquad
 \epsilon\in\{+1,-1\}.                               \tag{7}
\]
This valuation sign is unrelated to the physical sign $\sigma$.
For $\epsilon=-1$, clearing denominators gives
\[
 W^2=-(v^2+2)(v^4+v^2-4).
\]
The coprime parity cases give $3$, $6$, or $7$ modulo $8$, or an odd
$2$-adic valuation, so there is no rational point.  For $\epsilon=1$, the
discriminant of (6) is exactly (1), and (2) follows.

The zeros and poles divided out above are precisely those in
Table~\ref{tab:exceptions}.  Direct evaluation uses the original characters,
so even the fibre $L=M=0$ is not lost.
\end{proof}

\section{Square abscissas}

\begin{theorem}[Square-abscissa theorem]\label{thm:square}
The points of $D_+(\Q)$ with rational-square affine abscissa, together with
the point at infinity, are exactly
\[
 \infty,\qquad(0,0),\qquad(1,1),\qquad(1,-1).        \tag{8}
\]
\end{theorem}

\begin{proof}
Let $J=\operatorname{Jac}(D_+)$,
$D=[(1,1)-\infty]$, and $T=[(0,0)-\infty]$.
A complete Schaefer--Stoll descent in
\[
 \Q\times\Q[z]/(z^4-16z+20),\qquad S=\{2,5,17\},
\]
starts with $512$ square-norm classes and leaves
\[
 512\longrightarrow64\longrightarrow8\longrightarrow4.
\]
The last four are $0,D,T,D+T$.  The condition at $17$ is active.  Exact
good-reduction arithmetic proves
\[
 \operatorname{rank}J(\Q)=1,\qquad
 J(\Q)_{\rm tors}=\langle T\rangle\simeq\mathbf Z/2\mathbf Z,
\]
and $D,T$ are independent modulo $2$.

No primitive-generator assumption is used.  Write
$J(\Q)=\mathbf ZG\oplus\langle T\rangle$ and, after a torsion shift,
$D=mG$ with $m$ odd.  Exact Cantor arithmetic gives
\[
\begin{array}{c|cc}
p&\#J(\mathbf F_p)&\operatorname{ord}(\bar D)\\ \hline
7&42&21\\
23&462&231.
\end{array}
\]
If $r=\operatorname{ord}(\bar G)$ and $N$ is $21$ or $231$, then
$r/\gcd(r,m)=N$ and $r\mid2N$.  Oddness of $m$ excludes $r=2N$; hence
$r=N$ and $m$ is invertible modulo $N$.  The local $D$-coordinates of a
global class therefore agree modulo $21$.

At $7$, square abscissas give
\[
 (c,e)=(0,0),(\pm1,0),(\pm5,0),(0,1).
\]
At $23$, the torsion-zero coordinates reduce modulo $21$ to
\[
 c=0,\pm1,\pm2,\pm3,\pm4,\pm9,
\]
and the only torsion-one class is $(0,1)$.  Thus $\pm5$ are removed and only
the four disks containing the points in (8) survive.

Exact rational Cantor arithmetic for $42D$ gives normalized $7$-adic
logarithm direction $(1,4)$.  An annihilating differential therefore has
reduction
\[
 (4-a)\frac{da}{R}.
\]
It is nonzero in every surviving disk.  The only possible affine zero is
$a=4$, where the right side of (2) is the nonsquare $5$ modulo $7$; the
second basis differential is a unit at infinity.  Since $7>2g=4$, the
Coleman disk bound gives at most one rational point in each surviving disk.
The four points in (8) occupy them.  A second implementation independently
recomputes the finite Jacobian orders from exhaustive counts over
$\mathbf F_p$ and $\mathbf F_{p^2}$.
\end{proof}

\section{Exceptional lifts and conclusion}

\begin{table}[ht]
\centering
\small
\begin{tabular}{p{.20\textwidth}p{.31\textwidth}p{.39\textwidth}}
\toprule
Stage & Rational candidate & Exact disposition\\
\midrule
Compact inverse & $z=0,\pm1$ & $s=0$ or $s=1$.\\
Reciprocal charts & $S,K=0,\pm1$ and the negative radial zero &
Boundary; a constant-one quartic has no rational root; the radial
quadratic has discriminant $2^8\cdot5$.\\
$C_8\leftrightarrow E_1$ & $\mathcal U=0,\pm1,1/2$; $x=2$, $x^2=5$,
$\pm4P$, infinity & No rational inverse, boundary or wrong sign,
irreducible fibre, or a common character with odd $2$-adic valuation;
infinity has $\mathcal U\mathcal V=1/2$.\\
Character identity & $X=0,1,3,-1$ and zeros/poles of $g,L,M$ &
Direct values are boundary or nonsquare.  At $(X,Y)=(1,0)$,
$L=M=0$ but the undivided character is $1/2$.\\
Final map to $D_+$ & $\infty,(0,0),(1,\pm1)$ &
Infinity is boundary; $(0,0)$ would require $r^2=5$; the $a=1$ fibres
have character $1/2$ or $0$.\\
\bottomrule
\end{tabular}
\caption{Exact exceptional-fibre audit for Theorem~\ref{thm:reduction}.}
\label{tab:exceptions}
\end{table}

\begin{theorem}[Sharipov's second cuboid conjecture]
For coprime positive integers $p\ne q$, the polynomial $Q_{p,q}(t)$ is
irreducible over $\Q$.
\end{theorem}

\begin{proof}
By Theorem~\ref{thm:square}, a point arising from
Theorem~\ref{thm:reduction} must lie above one of the four values in (8).
Table~\ref{tab:exceptions} excludes every lift: in particular, at $a=1$ the
two translated elliptic points are $(1,0)$ and $(3,4)$, with undivided
common-character values $1/2$ and $0$.  Thus no $2+3$ factor exists.

Theorem~\ref{thm:linear} excludes $1+4$.  For the symmetric $5+5$ case, the
pinned source archive of arXiv:2601.04240 has SHA--256
\[
\texttt{96d7e4b7d3cb75af238a0ca8cd143a8d555daf6ac9e8c2cc94a099b5aea04760}.
\]
The replay extracts its complete second code block and exactly recomputes the
elimination polynomial, discriminant factorization, rational-root tests,
positive-root isolations, and all seven Sturm counts.  Every admissible
positive parameter interval has zero real candidate roots; the special
parameter $s=1$ is factored separately and is inadmissible.  Hence the
symmetric $5+5$ split is impossible.  Proposition~\ref{prop:factor} now
excludes every proper factorization.
\end{proof}

\section*{Reproducibility}

The proof-critical commands are:
\begin{verbatim}
PYTHON_BIN=.venv/bin/python bash scripts/reproduce.sh
PYTHON_BIN=.venv/bin/python bash scripts/replay_5plus5.sh
tectonic paper/main.tex
tectonic paper/clean.tex
\end{verbatim}
The reproduction script runs, in order, the compact model, reciprocal
quotient, elliptic transport, common-character and valuation descents,
genus-two 2-descent, two-prime sieve, formal logarithm, and independent
finite-Jacobian replay.  Their JSON outputs are regenerated and compared
exactly by the test suite.

The direct normalization claim C025, finite congruence sieves, Frobenius
fingerprints, genus-$33$ covers, and bounded-height searches are not
dependencies of this proof.

\end{document}
